%% kapa: exam=1 round=22 year=2011 period=2 form=A subject=appraisal_law subject_ko=부동산관계법규 seq=14 number=54
%% source: Q-Net 감정평가사 기출문제 (cst00309) artlSeq=5036790 file=extracted/1008/5036790/2011년 제22회 감정평가사 1차시험 기출문제/2교시.hwp
%% conversion: kordoc 4.24.51 -> tools/gen1.py (auto); review=proofread-gpt-astra; flags=none
\begin{question}{54}
국토의 계획 및 이용에 관한 법령상 토지거래계약허가제도에 대한 설명으로 옳지 않은 것은? (다툼이 있으면 판례에 의함)
\begin{choices}
  \item 토지거래허가구역의 지정은 토지의 투기적인 거래가 성행하거나 지가가 급격히 상승하는 지역 등에 대해 국토해양부장관이 행한다.
  \item 행정청의 토지거래허가구역의 지정행위는 개인의 권리 내지 법률상의 이익을 구체적으로 규제하는 효과를 가져오게 하는 행정청의 처분에 해당하지 않는다.
  \item 토지거래허가제는 사유재산제도를 부정하는 것이 아니라 토지의 투기적 거래억제를 위한 것으로 헌법상의 비례원칙이나 과잉금지원칙에 반하지 아니한다는 것이 헌법재판소의 입장이다.
  \item 허가를 받지 않고 체결한 계약의 효력에 대해 판례는 확정적인 무효와 유동적인 무효로 구별하고 있다.
  \item 벌칙적용대상인 허가없이 토지등의 거래계약을 체결한 행위라 함은 처음부터 허가를 배제하거나 잠탈하는 내용의 계약을 체결하는 것을 의미한다.
\end{choices}
\end{question}
\answer{2}
